\section{Modelo de adversario}
\label{sec:adversario}
Aqu\'i definimos contra qui\'en nos defendemos: qu\'e puede intentar cada
parte y qu\'e queda fuera del modelo.

\begin{table}[ht]
\centering
\small
\begin{tabularx}{\textwidth}{@{}l l X@{}}
\toprule
Actor & Se porta as\'i & Qu\'e intenta \\
\midrule
Servidor & Cumple el protocolo pero es curioso &
Sigue los pasos, pero intenta adivinar el $i$ mirando los mensajes
$(u,\{E_j\})$. No lo logra si el problema DH es dif\'icil. \\
Cliente & Quiere m\'as de lo permitido &
Cumple para pedir, pero intenta sacar $K_j$ de\'indices que no pidi\'o
($j\ne i$) con una sola respuesta. No puede sin romper el problema DH. \\
Red & Esp\'ia y modifica &
Mira el tr\'afico o cambia bytes. Cualquier cambio debe ser detectado
por el cifrado AEAD y el control de versi\'on. \\
\bottomrule
\end{tabularx}
\caption{Modelo de adversario del sistema.}
\end{table}

\textbf{Supuestos y l\'imites del modelo}:
(i)~PIR, extensi\'on de OT, OT adaptativo con OPRF ni bases de millones de
registros;
(ii)~clientes o servidores totalmente tramposos (solo cubrimos servidor
curioso que cumple el protocolo y cliente que pide de m\'as; si el servidor
manda datos falsos solo lo detectamos por integridad, como en el ejemplo de
los art\'iculos vac\'ios de la Secci\'on~11.6.4 del libro);
(iii)~pruebas de seguridad para cualquier combinaci\'on de programas
(componibilidad);
(iv)~ocultar IP, frecuencia, momento o volumen de consultas, ni pagos o
licenciamiento real.
